\chapter{The Hyperbolic Plane and the Modular Group}

\begin{center}\itshape
The hyperbolic plane is the most symmetric surface of negative curvature, and its symmetries are the matrices of determinant one. Points, directions,
and the motion along geodesics all become matrix multiplication. Inside this group sits $SL_2(\Z)$, which tiles the plane with triangles; a straight line
crossing the tiles writes down, letter by letter, the continued fraction of the point where it ends.
\end{center}

In 1880 Henri Poincaré, then twenty-six and teaching at Caen, was studying functions on the upper half-plane that are unchanged by a discrete group of
transformations $z\mapsto\frac{az+b}{cz+d}$, functions he named after Lazarus Fuchs. For two weeks he tried to prove that they could not exist, then found a family of
them, and then left on a geological excursion organized by the School of Mines. In \emph{Science et méthode} (1908) he described what happened in the
town of Coutances: the travels had made him forget his mathematics, and as he put his foot on the step of an omnibus the idea came to him, with nothing in
his previous thoughts seeming to have prepared it, that the transformations he had used to define his Fuchsian functions were identical with those of
non-Euclidean geometry. He did not check it then, but went on with a conversation already begun, ``with perfect certainty''; back in Caen he verified it at
leisure.

The non-Euclidean geometry of Bolyai and Lobachevsky had been a curiosity for half a century; Volume~8, Chapter~10 told that story and built the
half-plane model. Poincaré's observation turned it into a working tool: the groups of the theory of automorphic functions are groups of isometries of the
hyperbolic plane, and their fundamental domains are hyperbolic polygons. This chapter takes the next step. The full isometry group is $PSL_2(\R)$; the
unit tangent vectors of the plane \emph{are} the elements of this group, and following geodesics is multiplying on the right by a diagonal matrix. Inside sits the
modular group $SL_2(\Z)$ of Volume~5 and Volume~12, and Caroline Series' theorem of 1985 closes the chapter: continued fractions are the records of
geodesics crossing the Farey tessellation.

\begin{reminder}[from Volume~8, Chapter 10, and Chapter 4]
$\mathbb H=\{z=x+iy:y>0\}$ with length element $\frac{|\dd z|}y$; its curvature is $-1$ (Example~\ref{ex:halfplane}). The maps $z\mapsto\frac{az+b}{cz+d}$, $ad-bc=1$, and the reflection $z\mapsto-\bar z$ preserve lengths, distances and areas (Volume~8,
Proposition~J.4). The \emph{h-lines} are the vertical half-lines and the half-circles centred on $\R$; through two points passes exactly one, and an isometry above maps any
two points $z\ne w$ to $i$ and $it$ with $t>1$ (Lemma~10.5). The distance is given by
\[\cosh d(z,w)=1+\frac{|z-w|^2}{2\operatorname{Im}z\operatorname{Im}w},\]
and the arc of h-line between two points is the unique shortest path (Theorem~10.6). A triangle with angles $\alpha,\beta,\gamma$ has area $\pi-\alpha-\beta-\gamma$, ideal vertices counting as
angle $0$ (Theorem~10.8). The Cayley map $C(z)=\frac{z-i}{z+i}$ carries $\mathbb H$ to the unit disc, where the length element is $\frac{2|\dd w|}{1-|w|^2}$, the h-lines through $0$ are the diameters, and
$d(0,w)=2\operatorname{artanh}|w|$ (Proposition~10.10).
\end{reminder}

\section{The isometry group}

An \emph{isometry} of $\mathbb H$ is a map $f\colon\mathbb H\to\mathbb H$ with $d(f(z),f(w))=d(z,w)$. Since $g$ and $-g$ act in the same way, the maps $z\mapsto\frac{az+b}{cz+d}$ form the group
$PSL_2(\R)=SL_2(\R)/\{\pm I\}$.

\begin{lemma}\label{lem:hypcircles}
The hyperbolic circle $\{z:d(z,ib)=r\}$, $b>0$, is a Euclidean circle with centre $ib\cosh r$ on the imaginary axis. Two distinct such circles meet in at most two points,
which are symmetric with respect to the imaginary axis.
\end{lemma}

\begin{proof}
By the distance formula, $d(z,ib)=r$ means $|z-ib|^2=2by(\cosh r-1)$, that is $x^2+y^2-2by\cosh r+b^2=0$, or $|z-ib\cosh r|^2=b^2\sinh^2r$. Subtracting the equations of two such
circles with different $(b,r)$ gives a linear equation in $y$ alone (the $x^2+y^2$ cancel), or a contradiction; so the common points lie on a horizontal line, which meets each
circle in at most two points $\pm x+iy$.
\end{proof}

\begin{theorem}[the isometry group]\label{thm:isometries}
Every isometry of $\mathbb H$ is of the form $z\mapsto g(z)$ or $z\mapsto g(-\bar z)$ with $g\in PSL_2(\R)$. In particular the orientation-preserving isometries form the group $PSL_2(\R)$.
\end{theorem}

\begin{proof}
Let $f$ be an isometry. By the reminder, there is $g\in PSL_2(\R)$ with $g(f(i))=i$ and $g(f(2i))=it$ for some $t>1$; since $d(i,it)=\log t$ must equal $d(i,2i)=\log2$, $t=2$. So $h=g\circ f$ is
an isometry fixing $i$ and $2i$. For every $z$, the point $h(z)$ lies at distance $d(z,i)$ from $i$ and $d(z,2i)$ from $2i$, so on two distinct circles of Lemma~\ref{lem:hypcircles} through $z$:
hence $h(z)\in\{z,-\bar z\}$.

Suppose that $h(z_0)=z_0$ for some $z_0=x_0+iy_0$ with $x_0\ne0$. For $w=u+iv$ with $u\neq0$, the points $w$ and $-\bar w$ have the same imaginary part, and
$|w-z_0|^2-|-\bar w-z_0|^2=(u-x_0)^2-(u+x_0)^2=-4ux_0\neq0$, so they have different distances to $z_0$. Since $h(w)$ is at distance $d(w,z_0)$ from $h(z_0)=z_0$, $h(w)=w$. Thus either $h$ fixes
every point off the imaginary axis, and then every point, or it fixes none, and then $h(z)=-\bar z$ for all $z$ (on the axis the two choices agree). So $f=g^{-1}$ or $f(z)=g^{-1}(-\bar z)$.
\end{proof}

\begin{proposition}[geodesic polar coordinates]\label{prop:hyppolar}
For $r>0$ and $\vartheta\in\R$ let $P(r,\vartheta)=C^{-1}\bigl(\tanh\frac r2\,e^{i\vartheta}\bigr)$. Then $r=d(i,P(r,\vartheta))$, the curves $\vartheta=\mathrm{const}$ are the unit-speed geodesics from $i$, and the metric reads
\[\dd r^2+\sinh^2r\,\dd\vartheta^2.\]
The circle of radius $r$ has length $2\pi\sinh r$ and the disc of radius $r$ has area $2\pi(\cosh r-1)$.
\end{proposition}

\begin{proof}
Work in the disc, where $i$ becomes $0$, with $w=\tanh\frac r2e^{i\vartheta}$. The distance claim is $d(0,w)=2\operatorname{artanh}|w|$, and the rays from $0$ are the h-lines through $0$, traversed
at unit speed since $r$ is the distance. Writing $\rho=\tanh\frac r2$, $|\dd w|^2=\dd\rho^2+\rho^2\dd\vartheta^2$ with $\dd\rho=\frac12\cosh^{-2}\frac r2\dd r$, and $\frac2{1-\rho^2}=2\cosh^2\frac r2$. So the squared length
element is
\[4\cosh^4\tfrac r2\Bigl(\tfrac14\cosh^{-4}\tfrac r2\,\dd r^2+\tanh^2\tfrac r2\,\dd\vartheta^2\Bigr)=\dd r^2+\bigl(2\sinh\tfrac r2\cosh\tfrac r2\bigr)^2\dd\vartheta^2,\]
and $2\sinh\frac r2\cosh\frac r2=\sinh r$. The circle has length $\int_0^{2\pi}\sinh r\dd\vartheta$, and the area element is $\sinh r\dd r\dd\vartheta$.
\end{proof}

This keeps the promise of Exercise~\ref{ex:geodlength}: the metric of the half-plane in geodesic polar coordinates is $\dd r^2+\sinh^2r\,\dd\vartheta^2$, and with $G=\sinh^2r$ the formula
$K=-\frac{(\sqrt G)''}{\sqrt G}$ gives $K=-1$ once more. A circle of radius $10$ has length $2\pi\sinh10\approx69\,000$, nearly seven thousand times its radius: in the hyperbolic plane, room grows exponentially.

\section{Points, directions and matrices}

Let $K=SO(2)$, $A=\{a_t=\mathrm{diag}(e^{t/2},e^{-t/2}):t\in\R\}$ and $N=\{n_x=\bigl(\begin{smallmatrix}1&x\\0&1\end{smallmatrix}\bigr):x\in\R\}$. Write $k_\theta=\bigl(\begin{smallmatrix}\cos\theta&-\sin\theta\\\sin\theta&\cos\theta\end{smallmatrix}\bigr)$. The three kinds of
one-parameter subgroups of $SL_2(\R)$ met in Chapter~7 are, up to conjugacy, these three.

\begin{theorem}[Iwasawa decomposition]\label{thm:iwasawa}
\begin{enumerate}
\item[(a)] $SL_2(\R)$ acts transitively on $\mathbb H$, and the stabilizer of $i$ is $K$.
\item[(b)] Every $g\in SL_2(\R)$ is uniquely $g=n_xa_tk_\theta$ with $x,t\in\R$ and $\theta\in\R/2\pi\Z$, and $(x,t,\theta)$ depends smoothly on $g$.
\item[(c)] The differential of $k_\theta$ at $i$ is multiplication by $e^{-2i\theta}$: $k_\theta$ turns the tangent vectors at $i$ by the angle $-2\theta$.
\end{enumerate}
\end{theorem}

\begin{proof}
(a) $a_t(i)=e^ti$ and $n_x(z)=z+x$, so $n_xa_t$ maps $i$ to $x+ie^t$, any point of $\mathbb H$. A matrix fixes $i$ when $ai+b=i(ci+d)$, that is $a=d$, $b=-c$, and with
$a^2+b^2=1$ this is $K$. (b) Given $g$, let $g(i)=x+iy$ and $t=\log y$; then $(n_xa_t)^{-1}g$ fixes $i$, so it is some $k_\theta$. Uniqueness: if $n_xa_tk_\theta=n_{x'}a_{t'}k_{\theta'}$, apply both to $i$ to get $x=x'$,
$t=t'$, and then $k_\theta=k_{\theta'}$. Explicitly, for $g=\bigl(\begin{smallmatrix}a&b\\c&d\end{smallmatrix}\bigr)$, $g(i)=\frac{(ai+b)(d-ci)}{c^2+d^2}$, so $x=\frac{ac+bd}{c^2+d^2}$, $e^t=\frac1{c^2+d^2}$, and $k_\theta=a_{-t}n_{-x}g$: all smooth. (c) The
derivative of $z\mapsto\frac{az+b}{cz+d}$ is $(cz+d)^{-2}$; for $k_\theta$ at $i$ it is $(i\sin\theta+\cos\theta)^{-2}=e^{-2i\theta}$.
\end{proof}

The \emph{unit tangent bundle} $T^1\mathbb H$ is the set of pairs $(z,v)$ with $v$ a tangent vector at $z$ of hyperbolic length $1$, that is, Euclidean length $\operatorname{Im}z$. Let $v_0$ be the upward
unit vector at $i$. The \emph{geodesic flow} $\phi_t$ moves $(z,v)$ along the unit-speed geodesic $\gamma$ with $\gamma(0)=z$, $\gamma'(0)=v$: $\phi_t(z,v)=(\gamma(t),\gamma'(t))$. An isometry $g$ maps
unit-speed geodesics to unit-speed geodesics, because geodesic arcs are the unique shortest paths; so $g\circ\phi_t=\phi_t\circ dg$.

\begin{theorem}[the geodesic flow is a matrix product]\label{thm:geodflow}
The map $\Psi\colon PSL_2(\R)\to T^1\mathbb H$, $\Psi(g)=(g(i),dg_i(v_0))$, is a bijection, and it turns the geodesic flow into right multiplication by $a_t$:
\[\phi_t(\Psi(g))=\Psi(ga_t).\]
\end{theorem}

\begin{proof}
\emph{Bijection.} By Theorem~\ref{thm:iwasawa}(a), every point is $g(i)$; the elements of $PSL_2(\R)$ mapping $i$ to $z$ form a coset $gK$, and $dg_i\circ d(k_\theta)_i$ runs once over the unit vectors at $z$
as $\theta$ runs over $\R/\pi\Z$ (part (c); $k_{\theta+\pi}=-k_\theta$ is the same element of $PSL_2(\R)$).

\emph{Flow.} The curve $t\mapsto ie^t=a_t(i)$ is the vertical h-line through $i$, at unit speed since $\frac{|\dd y|}y=\dd t$; its velocity at time $t$ is the upward vector $ie^t$, which is $d(a_t)_i(v_0)$ because
$a_t$ is $z\mapsto e^tz$. So $\phi_t(\Psi(I))=\Psi(a_t)$. For general $g$, $\phi_t(\Psi(g))=\phi_t(dg(i,v_0))=dg\,\phi_t(i,v_0)=dg\,\Psi(a_t)=\Psi(ga_t)$.
\end{proof}

Similarly, right multiplication by $n_s$ moves $(z,v)$ along the horocycle tangent to $v$: for $g=I$, $n_s(i)=i+s$ runs along the horizontal line through $i$, a curve which stays at
constant height and is not a geodesic. Geometry has become algebra: the left action of $PSL_2(\R)$ moves the whole picture rigidly, the right action of $A$ is the geodesic flow, and
the right action of $N$ the horocycle flow. When a discrete group $\Gamma$ acts on the left, the flows pass to the quotient $\Gamma\backslash PSL_2(\R)$, the unit tangent bundle of the surface
$\Gamma\backslash\mathbb H$, and the dynamics of geodesics on that surface becomes the dynamics of a matrix product on a homogeneous space; Eberhard Hopf proved in 1936 that the
geodesic flow of a surface like the modular surface below is ergodic, and Volume~10 develops the ergodic theory behind such statements.

\begin{theorem}[elliptic, parabolic, hyperbolic]\label{thm:trace}
Let $g\in SL_2(\R)$, $g\neq\pm I$.
\begin{enumerate}
\item[(a)] If $|\operatorname{tr}g|<2$, $g$ has exactly one fixed point $z_0\in\mathbb H$ and is conjugate to some $k_\theta$: it is a rotation about $z_0$ (\emph{elliptic}).
\item[(b)] If $|\operatorname{tr}g|=2$, $g$ has no fixed point in $\mathbb H$, one on $\R\cup\{\infty\}$, and is conjugate to $\pm n_1$ or $\pm n_{-1}$ (\emph{parabolic}).
\item[(c)] If $|\operatorname{tr}g|>2$, $g$ is conjugate to $\pm\mathrm{diag}(\lambda,\lambda^{-1})$ with $\lambda>1$ (\emph{hyperbolic}). It fixes two points of $\R\cup\{\infty\}$, maps the h-line joining them, its
\emph{axis}, to itself, and translates it by $\ell=2\log\lambda$. Moreover $d(z,g(z))\ge\ell$ for every $z$, with equality exactly on the axis, and
\[2\cosh\frac\ell2=|\operatorname{tr}g|.\]
\end{enumerate}
\end{theorem}

\begin{proof}
Conjugation does not change the trace, and $hgh^{-1}$ fixes $h(z)$ when $g$ fixes $z$. The fixed points of $g=\bigl(\begin{smallmatrix}a&b\\c&d\end{smallmatrix}\bigr)$ solve $cz^2+(d-a)z-b=0$, of
discriminant $(d-a)^2+4bc=(a+d)^2-4$. (a) Here $c\ne0$ (otherwise $ad=1$ and $|a+d|\ge2$), and the discriminant is negative: one root $z_0$ in $\mathbb H$, its conjugate in the lower half-plane.
If $h(i)=z_0$, then $h^{-1}gh$ fixes $i$, so lies in $K$. (c) The eigenvalues are real, $\lambda^{\pm1}$ with $|\lambda|>1$, with real eigenvectors $v_1,v_2$; the matrix $P$ with columns $v_1$, $v_2$
diagonalizes $g$, and so does $P$ with $v_2$ replaced by $-v_2$, so we can take $\det P>0$ and then scale it to $\det P=1$. Replacing $g$ by $-g$ if needed, $g$ is conjugate to $z\mapsto\lambda^2z$, which fixes $0$ and
$\infty$ and translates the imaginary axis by $d(i,\lambda^2i)=2\log\lambda$. For $z=x+iy$,
\[\cosh d(z,\lambda^2z)=1+\frac{(\lambda^2-1)^2(x^2+y^2)}{2\lambda^2y^2}\ge1+\frac{(\lambda^2-1)^2}{2\lambda^2}=\frac{\lambda^2+\lambda^{-2}}2=\cosh(2\log\lambda),\]
with equality iff $x=0$. Finally $|\operatorname{tr}g|=\lambda+\lambda^{-1}=2\cosh\log\lambda=2\cosh\frac\ell2$. (b) The only eigenvalue is $\pm1$, and $g\ne\pm I$ is not diagonalizable, so in a suitable basis
$g=\pm\bigl(\begin{smallmatrix}1&u\\0&1\end{smallmatrix}\bigr)$ with $u\ne0$, again with a basis change of determinant $1$; conjugating by $a_t$ multiplies $u$ by $e^t$, which makes $u=\pm1$. The map $z\mapsto z\pm1$ has no fixed
point in $\mathbb H$, and fixes $\infty$ only.
\end{proof}

\section{The modular group}

The modular group $SL_2(\Z)$ acts on $\mathbb H$ through $PSL_2(\Z)$. Volume~5 proved that it is generated by $T=\bigl(\begin{smallmatrix}1&1\\0&1\end{smallmatrix}\bigr)$ and $S=\bigl(\begin{smallmatrix}0&-1\\1&0\end{smallmatrix}\bigr)$
(Volume~5, Proposition~12.1) and that every orbit meets
\[F=\bigl\{z\in\mathbb H:|\operatorname{Re}z|\le\tfrac12,\ |z|\ge1\bigr\}\]
(Volume~5, Theorem~12.2); Volume~12, Theorem~14.1 showed that two points of $F$ are in the same orbit only if they lie on the boundary and are glued by $T$ or $S$. So $F$ is a \emph{fundamental
domain}: its images tile $\mathbb H$ (Figure~\ref{fig:modular}). Geometrically, $F$ is a hyperbolic triangle: its sides are the h-lines $\operatorname{Re}z=\pm\frac12$ and the unit circle, its vertices are
$\rho=e^{2\pi i/3}$, $\rho+1=e^{i\pi/3}$ and the ideal vertex $\infty$. At $\rho$, the unit circle has a tangent making the angle $\frac\pi6$ with the vertical (the radius to $\rho$ makes the angle $\frac\pi6$ with
the horizontal), so the interior angle is $\frac\pi3$, and likewise at $\rho+1$. By Theorem~J.8 of Volume~8,
\[\operatorname{area}(F)=\pi-\frac\pi3-\frac\pi3-0=\frac\pi3.\]
The quotient $SL_2(\Z)\backslash\mathbb H$ is the \emph{modular surface}: the triangle $F$ with its vertical sides glued by $T$ and the two halves of its bottom arc glued by $S$. It is a sphere with
a cusp, of finite area $\frac\pi3$ although it is not compact, with two cone points at $i$ (angle $\pi$) and $\rho$ (angle $\frac{2\pi}3$). Its closed geodesics come from hyperbolic elements of
$SL_2(\Z)$: the axis of such an element projects to a closed curve of length $\ell$, and by Theorem~\ref{thm:trace} the length is computed from the trace, an integer $|\operatorname{tr}g|\ge3$
(Exercise~\ref{ex:shortest}).

\begin{figure}[ht]
\centering
\fig{ch09-modular}
\caption{The tessellation of $\mathbb H$ by the images of the fundamental domain $F$ (shaded) under $SL_2(\Z)$. The tiles accumulate at every rational point of the real axis; each is a
hyperbolic triangle of area $\frac\pi3$, and they are all congruent, however small some of them look.}
\label{fig:modular}
\end{figure}

\section{Farey tessellation and cutting sequences}

Write rational numbers in lowest terms $\frac pq$ with $q>0$, and $\infty=\frac10$. Two of them $\frac pq$, $\frac rs$ are \emph{Farey neighbours} if $|ps-qr|=1$, and the h-line joining them is then a
\emph{Farey edge}. The edges from $\infty$ are the vertical lines $\operatorname{Re}z=n$, $n\in\Z$, since $|1\cdot q-0\cdot p|=1$ means $q=1$. Let $\Delta_0$ be the open ideal triangle with vertices $0$, $1$,
$\infty$, that is $\{0<\operatorname{Re}z<1,\ |z-\frac12|>\frac12\}$.

\begin{proposition}[the Farey tessellation]\label{prop:farey}
\begin{enumerate}
\item[(a)] The group $GL_2(\Z)$ permutes the Farey edges, acting by $z\mapsto\frac{az+b}{cz+d}$ if the determinant is $1$ and by $z\mapsto\frac{a\bar z+b}{c\bar z+d}$ if it is $-1$; these maps are isometries
of $\mathbb H$, and $SL_2(\Z)$ is transitive on the Farey edges.
\item[(b)] Two distinct Farey edges do not cross.
\item[(c)] No Farey edge meets $\Delta_0$, nor any of its images $g(\Delta_0)$, $g\in SL_2(\Z)$. The two sides of the edge from $\frac pq$ to $\frac rs$ ($ps-qr=-1$) are the triangles with third vertex
$\frac{p+r}{q+s}$ and $\frac{p-r}{q-s}$.
\end{enumerate}
\end{proposition}

\begin{proof}
(a) A matrix of $GL_2(\Z)$ maps $\binom pq$ to an integer vector, which is primitive since the inverse matrix is integral, and maps the pair of vectors $\binom pq,\binom rs$ to a pair whose determinant
is multiplied by $\pm1$; so $|ps-qr|$ is preserved. A map with determinant $-1$ is $z\mapsto-\bar z$ followed by an element of $SL_2(\Z)$ (indeed $\frac{a\bar z+b}{c\bar z+d}=g(-\bar z)$ for the matrix
$g$ obtained by changing the signs of the first column), an isometry by the reminder. Given neighbours with $ps-qr=1$, the matrix $\bigl(\begin{smallmatrix}p&r\\q&s\end{smallmatrix}\bigr)\in SL_2(\Z)$
maps $\infty$ to $\frac pq$ and $0$ to $\frac rs$, so the edge $(0,\infty)$ to the given one; if $ps-qr=-1$, exchange the roles of the two endpoints.

(b) Two h-lines cross exactly when their endpoints interleave on $\R\cup\{\infty\}$. By (a), we may take one of the edges to be $(0,\infty)$; an edge crossing it joins $\frac pq<0<\frac rs$, with $p<0<r$
and $q,s\ge1$, so $|ps-qr|=|p|s+qr\ge2$.

(c) A geodesic meeting the open set $\Delta_0$ without crossing its three sides (which are Farey edges) stays in $\Delta_0$ (distinct h-lines that meet do cross), so its endpoints are among $0,1,\infty$ and
it is a side: absurd. Images of $\Delta_0$ are handled by (a). The matrix
$g=\bigl(\begin{smallmatrix}r&p\\s&q\end{smallmatrix}\bigr)$ has determinant $rq-ps=1$ and $g(0)=\frac pq$ and $g(\infty)=\frac rs$; the triangles $\Delta_0$ and $\Delta_0-1$ on the two sides of $(0,\infty)$ are mapped to the
triangles with third vertices $g(1)=\frac{p+r}{q+s}$ and $g(-1)=\frac{r-p}{s-q}$.
\end{proof}

The Farey triangles are bigger tiles than the images of $F$: $\Delta_0$ has area $\pi$, three times that of $F$, and is made of three pieces of the modular tessellation. We only need (b) and (c).

Let $\gamma$ be an oriented h-line from $y$ to $x$, two distinct real numbers. Of the two open arcs of $\R\cup\{\infty\}$ cut out by $y$ and $x$, call \emph{right} of $\gamma$ the one swept by moving from $y$ to $x$
along the real axis without passing through $\infty$ when $y<x$ (that is the interval $(y,x)$), and through $\infty$ when $y>x$; the other arc is the \emph{left}. (Walking along $\gamma$, the right arc is on
your right hand.) When $\gamma$ crosses a Farey edge $e$ into a Farey triangle and leaves it through a second side $e'$, the sides $e$ and $e'$ share a vertex $v$; write $L$ if $v$ is on the left of
$\gamma$ and $R$ if it is on the right. The word in $L$ and $R$ so produced is the \emph{cutting sequence} of $\gamma$.

\begin{theorem}[Series, 1985]\label{thm:cutting}
Let $x>1$ be irrational with continued fraction $x=[a_0;a_1,a_2,\dots]$, and $-1<y<0$. After it crosses the edge $(0,\infty)$, the h-line from $y$ to $x$ has cutting sequence
\[L^{a_0}R^{a_1}L^{a_2}R^{a_3}\cdots\]
\end{theorem}

\begin{proof}
For $0\le k\le a_0$, the endpoints satisfy $y<k<x$, so $\gamma$ crosses the vertical edge $(k,\infty)$. Between the crossings of $(k,\infty)$ and $(k+1,\infty)$, for $k<a_0$, it traverses the Farey triangle
$\Delta_0+k$ (which contains no other edge, by Proposition~\ref{prop:farey}(c)); the shared vertex is $\infty$, on the left of $\gamma$ since $y<x$. That gives $L^{a_0}$. After crossing $(a_0,\infty)$, $\gamma$ is in
$\Delta_0+a_0$, with vertices $a_0$, $a_0+1$, $\infty$; since $y<a_0<x<a_0+1$ it leaves through the side $(a_0,a_0+1)$, and the next letter, about this triangle, is not $L$ (the shared vertex $a_0$ lies
in $(y,x)$, on the right).

To read the rest, apply $\psi(z)=\frac1{\bar z-a_0}$, an isometry preserving the Farey edges by Proposition~\ref{prop:farey}(a) (matrix $\bigl(\begin{smallmatrix}0&1\\1&-a_0\end{smallmatrix}\bigr)$, determinant $-1$). It
maps $a_0\mapsto\infty$, $\infty\mapsto0$, $a_0+1\mapsto1$, so the edge $(a_0,\infty)$ to $(\infty,0)$ and the triangle $\Delta_0+a_0$ to $\Delta_0$, and it maps $\gamma$ to the h-line $\gamma'$ from
$y'=\frac1{y-a_0}\in(-1,0)$ (as $y-a_0<-1$) to $x'=\frac1{x-a_0}=[a_1;a_2,\dots]>1$. On the real line, $\psi$ maps $(y,a_0)$ onto $(-\infty,y')$ and $(a_0,x)$ onto $(x',+\infty)$: it sends the right arc $(y,x)$ of
$\gamma$ onto the left arc of $\gamma'$, and so exchanges left and right. Therefore the cutting sequence of $\gamma$ after the crossing of $(a_0,\infty)$ is that of $\gamma'$ after the crossing of $(0,\infty)$,
with $L$ and $R$ exchanged. By what we proved, the latter begins with $L^{a_1}$ followed by a letter that is not $L$; so $\gamma$ continues with $R^{a_1}$, and induction gives the theorem.
\end{proof}

\begin{figure}[ht]
\centering
\fig{ch09-farey}
\caption{The Farey tessellation (edges between Farey neighbours, here with denominators up to $12$) and the h-line from $1-\sqrt2$ to $1+\sqrt2=[2;2,2,\dots]$. After crossing the edge from $0$ to $\infty$ it
passes two triangles with the vertex $\infty$ on its left, then two with a vertex on its right, and so on: the cutting sequence $LLRRLLRR\cdots$.}
\label{fig:farey}
\end{figure}

\begin{beauty}
Euclid's algorithm, run on a real number, is a walk in the hyperbolic plane. Each step $x\mapsto\frac1{x-\lfloor x\rfloor}$ of the continued fraction, the Gauss map of Volume~1, is the isometry $\psi$ that
moves the observer to the next block of triangles; quadratic irrationals, whose continued fractions are eventually periodic (Volume~1, Theorem~8.5), are the endpoints of geodesics that close up on
the modular surface; and the best rational approximations $\frac{p_n}{q_n}$ of $x$ are the vertices where the cutting sequence turns. Emil Artin used exactly this dictionary in 1924 to exhibit a
geodesic on the modular surface that comes arbitrarily close to every point in every direction: take an $x$ whose continued fraction contains every finite string of positive integers.
\end{beauty}

\begin{remark}
Series proved more: coding geodesics on the modular surface by their cutting sequences conjugates the geodesic flow, sampled at its crossings of the Farey edges, to the shift on sequences, and
the Gauss map is its one-sided version. The Gauss measure $\frac{\dd x}{(1+x)\log2}$ of the continued fraction map is a shadow of the invariant (Haar) measure of the geodesic flow on
$SL_2(\Z)\backslash SL_2(\R)$, which is how Gauss's 1812 problem on the digits of continued fractions connects with the ergodic theory of Volume~10.
\end{remark}

\section{Exercises}

\begin{exercise}[$\star$]\label{ex:iwasawaexplicit}
(a) Compute the Iwasawa decomposition of $\bigl(\begin{smallmatrix}2&1\\1&1\end{smallmatrix}\bigr)$. (b) Show that the geodesics through $i$ are the curves $t\mapsto k_\theta a_tk_\theta^{-1}(i)$, and that the
horocycles through $i$ are the curves $s\mapsto k_\theta n_sk_\theta^{-1}(i)$. What are they in the disc model?
\end{exercise}

\begin{exercise}[$\star\star$]\label{ex:shortest}
(a) Show that every hyperbolic element of $SL_2(\Z)$ has $|\operatorname{tr}g|\ge3$ and translation length $\ell\ge2\operatorname{arccosh}\frac32=4\log\varphi\approx1.925$, where $\varphi=\frac{1+\sqrt5}2$.
(b) Show that $g=\bigl(\begin{smallmatrix}2&1\\1&1\end{smallmatrix}\bigr)$ attains it, and find its axis. (c) What is the cutting sequence of the axis, oriented towards $\varphi$?
\end{exercise}

\begin{exercise}[$\star$]\label{ex:sqrt2axis}
Find $g\in SL_2(\Z)$ whose axis is the h-line of Figure~\ref{fig:farey}, from $1-\sqrt2$ to $1+\sqrt2$, and its translation length. Check that $g$ maps the pattern of triangles crossed by the axis to
itself, shifted by one period $LLRR$.
\end{exercise}

\begin{exercise}[$\star\star$]\label{ex:formsreduction}
To a positive definite binary quadratic form $Q=ax^2+bxy+cy^2$ ($a>0$, $D=b^2-4ac<0$) attach the root $\tau_Q=\frac{-b+\sqrt D}{2a}\in\mathbb H$ of $Q(\tau,1)=0$. (a) Show that for $g\in SL_2(\Z)$, the
form $Q\circ g$ (substituting $(x,y)\mapsto g(x,y)$) has root $g^{-1}(\tau_Q)$. (b) Show that $\tau_Q\in F$ if and only if $|b|\le a\le c$, Gauss's condition for $Q$ to be \emph{reduced}. (c) List the reduced forms of
discriminant $-20$ and $-23$. (The class numbers $2$ and $3$ count the orbits; compare Volume~12, Chapter~3.)
\end{exercise}

\begin{exercise}[$\star$]\label{ex:klein336}
(a) Show that the hyperbolic triangle with angles $\frac\pi2$, $\frac\pi3$, $\frac\pi7$ has area $\frac\pi{42}$. (b) A compact surface of genus $3$ with curvature $-1$ has area $8\pi$ (Corollary~\ref{cor:gaussbonnet}). How many
such triangles tile Klein's quartic curve, which is such a surface? (Hurwitz proved in 1893 that a surface of genus $g\ge2$ has at most $84(g-1)$ conformal symmetries; this count explains the number
$168$ of Klein's quartic.)
\end{exercise}

\begin{exercise}[$\star\star$]\label{ex:horocyclearea}
Let $B_h=\{z:\operatorname{Im}z>h\}$, a \emph{horoball}. (a) Show that the part of $F$ inside $B_h$ ($h\ge1$) has area $\frac1h$. (b) Show that $d(i+s,i)=2\operatorname{arsinh}\frac{|s|}2$: points on a horocycle drift apart
only logarithmically in the horocyclic length $|s|$, while points on a circle of radius $r$ are separated by up to $2r$ along an arc of length about $\pi\sinh r$.
\end{exercise}

\section{Hints}

\begin{hint}{ex:iwasawaexplicit}
Use the formulas in the proof of Theorem~\ref{thm:iwasawa}. For (b), conjugating by $k_\theta$ rotates directions at $i$.
\end{hint}

\begin{hint}{ex:shortest}
$|\operatorname{tr}g|$ is an integer, and $2\cosh\frac\ell2=|\operatorname{tr}g|$. The fixed points solve $z^2-z-1=0$.
\end{hint}

\begin{hint}{ex:sqrt2axis}
Solve $cz^2+(d-a)z-b=0$ with root $1+\sqrt2$, that is, proportional to $z^2-2z-1$.
\end{hint}

\begin{hint}{ex:formsreduction}
$Q(x,y)=a(x-\tau_Qy)(x-\bar\tau_Qy)$. For (b), $|\tau|^2=\frac ca$ and $\operatorname{Re}\tau=-\frac b{2a}$.
\end{hint}

\begin{hint}{ex:klein336}
Theorem~J.8 of Volume~8, and divide.
\end{hint}

\begin{hint}{ex:horocyclearea}
The area element is $\frac{\dd x\dd y}{y^2}$, and use the distance formula.
\end{hint}

\section{Solutions}

\begin{solution}{ex:iwasawaexplicit}
\idea The $N$ and $A$ parts locate the point $g(i)$; the $K$ part is the direction.
\writeup (a) With $a,b,c,d=2,1,1,1$: $c^2+d^2=2$, $x=\frac{ac+bd}2=\frac32$, $e^t=\frac12$. Then $k_\theta=a_{-t}n_{-x}g$: $n_{-3/2}g=\bigl(\begin{smallmatrix}1/2&-1/2\\1&1\end{smallmatrix}\bigr)$, and $a_{-t}=\mathrm{diag}(\sqrt2,\frac1{\sqrt2})$ gives
$\bigl(\begin{smallmatrix}1/\sqrt2&-1/\sqrt2\\1/\sqrt2&1/\sqrt2\end{smallmatrix}\bigr)=k_{\pi/4}$. So $g=n_{3/2}a_{-\log2}k_{\pi/4}$; check: $g(i)=\frac{2i+1}{i+1}=\frac32+\frac i2$. (b) $k_\theta$ fixes $i$ and rotates the directions there;
$k_\theta a_tk_\theta^{-1}(i)=k_\theta(ie^t)$ is the image of the vertical geodesic under a rotation about $i$, and all geodesics through $i$ arise this way (Theorem~\ref{thm:iwasawa}(c)). Likewise $k_\theta n_sk_\theta^{-1}(i)=k_\theta(i+s)$
rotates the horizontal line $\operatorname{Im}z=1$, a horocycle, about $i$. In the disc model, the geodesics through $0$ are the diameters, and the horocycle $\operatorname{Im}z=1$ (a circle through
$\infty$ tangent to $\R$ there) becomes the circle through $C(\infty)=1$ tangent to the unit circle and passing through $C(i)=0$, the circle $|w-\frac12|=\frac12$; the horocycles through $0$ are its
rotations.
\end{solution}

\begin{solution}{ex:shortest}
\idea Integer traces force a minimal translation, and the matrix $\bigl(\begin{smallmatrix}2&1\\1&1\end{smallmatrix}\bigr)$ is the square of the Fibonacci matrix.
\writeup (a) For $g\in SL_2(\Z)$ hyperbolic, $|\operatorname{tr}g|$ is an integer $>2$, so at least $3$, and $\ell=2\operatorname{arccosh}\frac{|\operatorname{tr}g|}2\ge2\operatorname{arccosh}\frac32$. Since $\cosh u=\frac32$ means $e^u=\frac{3+\sqrt5}2=\varphi^2$,
this is $4\log\varphi\approx1.9248$. (b) $\operatorname{tr}g=3$. Its fixed points solve $z^2+(1-2)z-1=0$, that is $z^2-z-1=0$: $\varphi$ and $-\frac1\varphi=1-\varphi$. The axis is the half-circle from $-\frac1\varphi\approx-0.618$ to
$\varphi\approx1.618$, and $g$ translates along it towards $\varphi$ (the attracting fixed point: $g^n(z)\to\varphi$, as ratios of consecutive Fibonacci numbers do, since $g=\bigl(\begin{smallmatrix}1&1\\1&0\end{smallmatrix}\bigr)^2$).
(c) $\varphi=[1;1,1,\dots]$ and $-\frac1\varphi\in(-1,0)$, so by Theorem~\ref{thm:cutting} the cutting sequence is $LRLRLR\cdots$; the period $LR$ corresponds to one application of $g$.
\end{solution}

\begin{solution}{ex:sqrt2axis}
\idea The axis of $g$ is the h-line joining the roots of $cz^2+(d-a)z-b$; we need that polynomial proportional to $z^2-2z-1$, with $ad-bc=1$.
\writeup Take $c=k$, $d-a=-2k$, $b=k$. Then $ad-bc=a(a-2k)-k^2=1$, that is $(a-k)^2=1+2k^2$. The smallest solution is $k=2$, $a-k=3$: $a=5$, $d=1$, $b=c=2$, so $g=\bigl(\begin{smallmatrix}5&2\\2&1\end{smallmatrix}\bigr)$, of trace
$6$, with $\lambda=3+2\sqrt2=(1+\sqrt2)^2$ and $\ell=2\log\lambda=4\log(1+\sqrt2)\approx3.525$. (The Pell equation $(a-k)^2-2k^2=1$ of Volume~1 appears; its solution $3+2\sqrt2$ is the unit of $\Z[\sqrt2]$ of
norm $1$.) Since $1+\sqrt2=[2;2,2,\dots]$, the cutting sequence is $LLRRLLRR\cdots$ (Theorem~\ref{thm:cutting}). Now $g=M^2$ with $M=\bigl(\begin{smallmatrix}2&1\\1&0\end{smallmatrix}\bigr)$, of determinant $-1$, which acts by
$z\mapsto2+\frac1{\bar z}$ (Proposition~\ref{prop:farey}(a)): the inverse of the map $\psi$ of the proof of Theorem~\ref{thm:cutting} with $a_0=2$. So $M$ maps the axis to itself, preserves the tessellation, and
shifts the crossed triangles by one block of two letters while exchanging $L$ and $R$; $g=M^2$ shifts them by the period $LLRR$.
\end{solution}

\begin{solution}{ex:formsreduction}
\idea A positive definite form is a point of $\mathbb H$ up to scaling, and $SL_2(\Z)$ acts on both in the same way; reduction is moving the point into $F$.
\writeup (a) $Q(x,y)=a(x-\tau y)(x-\bar\tau y)$ with $\tau=\tau_Q$. With $g=\bigl(\begin{smallmatrix}p&q\\r&s\end{smallmatrix}\bigr)$, $Q(px+qy,rx+sy)=a\bigl((p-\tau r)x+(q-\tau s)y\bigr)(\cdots)$, whose root in $x/y$ is
$\frac{\tau s-q}{p-\tau r}=g^{-1}(\tau)$, lying in $\mathbb H$. (b) $\operatorname{Re}\tau=-\frac b{2a}$ and $|\tau|^2=\tau\bar\tau=\frac ca$. So $|\operatorname{Re}\tau|\le\frac12$ iff $|b|\le a$, and $|\tau|\ge1$ iff $c\ge a$. (c) For $D=-20$: $b$ is even,
$|b|\le a\le c$, $ac=\frac{b^2+20}4$, and $a\le\sqrt{|D|/3}<2.6$. $b=0$: $ac=5$, $(1,0,5)$; $b=\pm2$: $ac=6$, $a=2$, $c=3$: $(2,\pm2,3)$, and the two are equivalent (they correspond to boundary points
glued by $T$). So two classes, $x^2+5y^2$ and $2x^2+2xy+3y^2$. For $D=-23$: $b$ odd, $a\le2$; $b=\pm1$, $ac=6$: $(1,1,6)$, $(2,\pm1,3)$; $(1,-1,6)$ is equivalent to $(1,1,6)$ (boundary), while $(2,1,3)$ and $(2,-1,3)$
correspond to interior points of $F$, not equivalent. Three classes.
\end{solution}

\begin{solution}{ex:klein336}
\idea Area is additive and Gauss--Bonnet fixes the total.
\writeup (a) $\pi-\frac\pi2-\frac\pi3-\frac\pi7=\pi\frac{42-21-14-6}{42}=\frac\pi{42}$. (b) $\int K\dd A=-\operatorname{area}=2\pi\chi=2\pi(2-2g)=-8\pi$ for $g=3$; so $8\pi/\frac\pi{42}=336$ triangles. Half of them are mirror images of the
others, and the $168$ orientation-preserving symmetries of the quartic permute the $168$ triangles of each kind simply transitively: $168=84(3-1)$, Hurwitz's bound attained.
\end{solution}

\begin{solution}{ex:horocyclearea}
\idea Near the cusp the modular surface is a thin horn: its area beyond height $h$ is $\frac1h$, finite although the horn is infinitely long.
\writeup (a) For $h\ge1$ the part of $F$ above height $h$ is $|x|\le\frac12$, $y>h$, of area $\int_{-1/2}^{1/2}\int_h^\infty\frac{\dd y}{y^2}\dd x=\frac1h$. (b) $\cosh d=1+\frac{s^2}2$, and $\cosh2u=1+2\sinh^2u$ gives $\sinh\frac d2=\frac{|s|}2$. So
$d=2\operatorname{arsinh}\frac{|s|}2\approx2\log|s|$ for large $s$: along a horocycle the distance grows logarithmically, whereas on a circle of radius $r$ two antipodal points are at distance $2r$ while the
arc between them has length $\pi\sinh r\approx\frac\pi2e^r$.
\end{solution}

\section*{Sources and further reading}
\addcontentsline{toc}{section}{Sources and further reading}

H.~Poincaré, \emph{Science et méthode} (Flammarion, 1908), Book~I, Chapter~3, ``L'invention mathématique''; English translation by G.~B.~Halsted in \emph{The Foundations of Science} (Science
Press, 1913). J.~Gray, \emph{Henri Poincaré: A Scientific Biography} (Princeton University Press, 2013), on the Fuchsian functions of 1880--81 and the Coutances episode.

C.~Series, ``The modular surface and continued fractions'', \emph{J.\ London Math.\ Soc.} (2) 31 (1985), 69--80; E.~Artin, ``Ein mechanisches System mit quasiergodischen Bahnen'', \emph{Abh.\ Math.\
Sem.\ Univ.\ Hamburg} 3 (1924), 170--175; E.~Hopf, ``Fuchsian groups and ergodic theory'', \emph{Trans.\ Amer.\ Math.\ Soc.} 39 (1936), 299--314; A.~Hurwitz, ``\"Uber algebraische Gebilde mit
eindeutigen Transformationen in sich'', \emph{Math.\ Ann.} 41 (1893), 403--442.

The treatment follows S.~Katok, \emph{Fuchsian Groups} (University of Chicago Press, 1992), Chapters~1--4; M.~Einsiedler and T.~Ward, \emph{Ergodic Theory with a view towards Number Theory}
(Springer, 2011), Chapter~9, for the geodesic flow as a matrix product and the link with continued fractions; and C.~Series, ``Continued fractions and hyperbolic geometry'', lecture notes of
the London Mathematical Society summer school, Loughborough, 2015, which draws the Farey pictures beautifully. For Klein's quartic, S.~Levy (ed.), \emph{The Eightfold Way} (MSRI
Publications 35, Cambridge University Press, 1999).
